% Generated by tools/edn_latex.py from reports/edn.md.
% Do not edit: change reports/edn.md and run `make edn`.
\ednentry{
  id = {EDN-30},
  title = {Pipeline configuration lives in \texorpdfstring{\texttt{params.\allowbreak{}yaml}}{params.yaml}, not in \texorpdfstring{\texttt{config.\allowbreak{}py}}{config.py}},
  date = {2026-09-29},
  milestone = {M2: Reproducibility},
  topic = {Project Structure, Reproducibility},
  participants = {@lukas2510},
  decision = {Everything configurable about the pipeline -- the seed, the age reference date, the scope bounds, the split ratios, the feature sets, the model variants and the SC-01 to SC-06 thresholds -- lives in \texttt{params.\allowbreak{}yaml} and is declared per stage under \texttt{params:} in \texttt{dvc.\allowbreak{}yaml}. Paths stay in \texttt{recommenditos/\allowbreak{}config.\allowbreak{}py}.},
  rationale = {\emph{Alternatives considered:}
\begin{itemize}[leftmargin=*, topsep=2pt]
\item \textbf{Option A (chosen): a \texttt{params.\allowbreak{}yaml} declared per stage.}
\newline Pros: \texttt{dvc repro} reruns exactly the stages a change affects, because DVC hashes the declared keys; \texttt{dvc params diff} shows what changed between two commits, which is evidence the report can cite; \texttt{dvc exp} becomes usable without a code change; and a value used by two stages cannot drift, because there is one copy of it.
\newline Cons: one more file to keep in step with the code, and a parameter typo is caught at run time rather than by the interpreter.
\item \textbf{Option B: Python constants in \texttt{config.\allowbreak{}py}, as the course demo does.}
\newline Pros: it is exactly what the demo shows, so it needs no defending; typos are import errors; no YAML parsing.
\newline Cons: DVC cannot see a constant, so a changed hyperparameter reruns nothing and \texttt{dvc repro} reports the pipeline as up to date while the code says otherwise -- which is the opposite of what the milestone is graded on. \texttt{dvc exp} cannot sweep it, and the demo's own \texttt{config.\allowbreak{}py} also reads an absolute \texttt{ROOT} from \texttt{.\allowbreak{}env} and raises at import without one.
\item \textbf{Option C: both, with \texttt{config.\allowbreak{}py} reading \texttt{params.\allowbreak{}yaml}.}
\newline Pros: one import surface for the stages.
\newline Cons: the indirection hides which stage reads which key, which is precisely the information \texttt{dvc.\allowbreak{}yaml} needs in order to rerun the minimum.
\end{itemize}
\par
\emph{Why:} The whole point of M2 is that a change reruns what it should and nothing else. Option B cannot deliver that for anything except code, and parameters are where most changes will happen once the stage tickets land. Paths were deliberately left out of \texttt{params.\allowbreak{}yaml}: they are not experiment variables, and deriving them from \texttt{\_\allowbreak{}\_\allowbreak{}file\_\allowbreak{}\_\allowbreak{}} keeps the project free of machine-specific absolute paths.},
  ai = {Alternative generation, Alternative assessment, Recommendation, Solution generation},
  response = {Accepted},
  assessment = {AI read the demo repository's \texttt{dvc.\allowbreak{}yaml} and \texttt{config.\allowbreak{}py} directly rather than from our notes, reported that the demo declares neither \texttt{params} nor \texttt{metrics} on any stage, and measured the consequence in a throwaway DVC repository: with a coarse \texttt{params:} declaration a single changed hyperparameter retrained every variant, and with a per-item declaration it retrained one. Lukas reviewed the measurement and accepted the recommendation.},
  aievidence = {Claude Code session on 2026-09-29 while building the pipeline skeleton: AI built a scratch DVC repository, ran \texttt{dvc repro} and \texttt{dvc status} under both declaration styles, and pasted the outputs showing one stage rerunning instead of three.},
  evidence = {\href{https://github.com/mlops-2627q1-mds-upc/MLOPS_Recommenditos/blob/main/params.yaml}{\texttt{params.\allowbreak{}yaml}}; \href{https://github.com/mlops-2627q1-mds-upc/MLOPS_Recommenditos/blob/main/dvc.yaml}{\texttt{dvc.\allowbreak{}yaml}}; \href{https://github.com/mlops-2627q1-mds-upc/MLOPS_Recommenditos/issues/32}{issue \#32}.},
}
